% Options for packages loaded elsewhere
\PassOptionsToPackage{unicode}{hyperref}
\PassOptionsToPackage{hyphens}{url}
\documentclass[
  ignorenonframetext,
]{beamer}
\newif\ifbibliography
\usepackage{pgfpages}
\setbeamertemplate{caption}[numbered]
\setbeamertemplate{caption label separator}{: }
\setbeamercolor{caption name}{fg=normal text.fg}
\beamertemplatenavigationsymbolsempty
% remove section numbering
\setbeamertemplate{part page}{
  \centering
  \begin{beamercolorbox}[sep=16pt,center]{part title}
    \usebeamerfont{part title}\insertpart\par
  \end{beamercolorbox}
}
\setbeamertemplate{section page}{
  \centering
  \begin{beamercolorbox}[sep=12pt,center]{section title}
    \usebeamerfont{section title}\insertsection\par
  \end{beamercolorbox}
}
\setbeamertemplate{subsection page}{
  \centering
  \begin{beamercolorbox}[sep=8pt,center]{subsection title}
    \usebeamerfont{subsection title}\insertsubsection\par
  \end{beamercolorbox}
}
% Prevent slide breaks in the middle of a paragraph
\widowpenalties 1 10000
\raggedbottom
\AtBeginPart{
  \frame{\partpage}
}
\AtBeginSection{
  \ifbibliography
  \else
    \frame{\sectionpage}
  \fi
}
\AtBeginSubsection{
  \frame{\subsectionpage}
}
\usepackage{iftex}
\ifPDFTeX
  \usepackage[T1]{fontenc}
  \usepackage[utf8]{inputenc}
  \usepackage{textcomp} % provide euro and other symbols
\else % if luatex or xetex
  \usepackage{unicode-math} % this also loads fontspec
  \defaultfontfeatures{Scale=MatchLowercase}
  \defaultfontfeatures[\rmfamily]{Ligatures=TeX,Scale=1}
\fi
\usepackage{lmodern}
\ifPDFTeX\else
  % xetex/luatex font selection
\fi
% Use upquote if available, for straight quotes in verbatim environments
\IfFileExists{upquote.sty}{\usepackage{upquote}}{}
\IfFileExists{microtype.sty}{% use microtype if available
  \usepackage[]{microtype}
  \UseMicrotypeSet[protrusion]{basicmath} % disable protrusion for tt fonts
}{}
\makeatletter
\@ifundefined{KOMAClassName}{% if non-KOMA class
  \IfFileExists{parskip.sty}{%
    \usepackage{parskip}
  }{% else
    \setlength{\parindent}{0pt}
    \setlength{\parskip}{6pt plus 2pt minus 1pt}}
}{% if KOMA class
  \KOMAoptions{parskip=half}}
\makeatother
\usepackage{longtable,booktabs,array}
\newcounter{none} % for unnumbered tables
\usepackage{calc} % for calculating minipage widths
\usepackage{caption}
% Make caption package work with longtable
\makeatletter
\def\fnum@table{\tablename~\thetable}
\makeatother
\setlength{\emergencystretch}{3em} % prevent overfull lines
\providecommand{\tightlist}{%
  \setlength{\itemsep}{0pt}\setlength{\parskip}{0pt}}
\usepackage{bookmark}
\IfFileExists{xurl.sty}{\usepackage{xurl}}{} % add URL line breaks if available
\urlstyle{same}
\hypersetup{
  hidelinks,
  pdfcreator={LaTeX via pandoc}}

\author{\texorpdfstring{}{}}
\date{}

\begin{document}

\section{Methodology: The CIF-AD-OODA Integration
Model}\label{sec:methodology}

\begin{frame}[allowframebreaks]{Methodology: The CIF-AD-OODA Integration
Model}
This section establishes the analytical framework used to evaluate
cognitive integrity threats across the ten critical domains examined in
this paper. We integrate three complementary theoretical
frameworks---the Cognitive Integrity Framework (CIF)
\cite{friedman2026cogsec1}, Axiomatic Design (AD)
\cite{suh1990principles, suh2001axiomatic}, and the OODA Loop
\cite{boyd1987patterns}---into a unified model for analyzing Goal
Hijacking attacks and their defenses.
\end{frame}

\begin{frame}[allowframebreaks]{Framework Integration}
\protect\phantomsection\label{framework-integration}
Each framework contributes a distinct analytical dimension:

\textbf{The Cognitive Integrity Framework (CIF)} provides the defense
mechanism vocabulary. Developed formally in Paper 1 of this series
\cite{friedman2026cogsec1} and validated computationally in Paper 2
\cite{friedman2026cogsec2}, CIF defines five canonical defense
mechanisms that protect agent cognitive states
\(\sigma_i = \langle \mathcal{B}_i, \mathcal{G}_i, \mathcal{I}_i, \mathcal{H}_i \rangle\)
against adversarial manipulation. These mechanisms operate at the
architectural, runtime, and coordination layers to maintain belief
consistency, goal alignment, and provenance verifiability.

\textbf{Axiomatic Design (AD)} provides the structural model. Suh's
Independence Axiom \cite{suh2001axiomatic} states that a good design
maintains the independence of Functional Requirements (FRs): solving for
\(\text{FR}_1\) must not compromise \(\text{FR}_2\). We represent this
as the Design Matrix equation:

where \(\{FR\}\) is the vector of Functional Requirements, \(\{DP\}\) is
the vector of Design Parameters, and \([A]\) is the Design Matrix. In an
uncoupled (safe) design, \([A]\) is diagonal. Goal Hijacking attacks
introduce off-diagonal terms, coupling previously independent
requirements. \begin{equation}
\{FR\} = [A]\{DP\}
\label{eq:meth_design_matrix}
\end{equation}

\textbf{The OODA Loop} provides the temporal model. Boyd's
Observe-Orient-Decide-Act cycle
\cite{boyd1987patterns, osinga2007science} captures the decision-making
process of autonomous agents. Goal Hijacking targets the \textbf{Orient}
phase---the synthesis stage where agents integrate observations with
prior knowledge, training data, and system prompts. By corrupting
Orientation, adversaries redirect the downstream Decide and Act phases
while preserving the illusion of rational behavior.

The integration insight is that Goal Hijacking operates simultaneously
across all three dimensions: it introduces transient coupling in the
Design Matrix (AD), corrupts the Orient phase (OODA), and is detectable
and defensible through the canonical mechanisms (CIF).
\end{frame}

\begin{frame}[allowframebreaks]{CIF Defense Mechanisms}
\protect\phantomsection\label{cif-defense-mechanisms}
Paper 1 \cite{friedman2026cogsec1} defines five canonical defense
mechanisms. We reproduce them here for reader convenience, as they form
the defense vocabulary applied across all ten domains:

{\def\LTcaptype{none} % do not increment counter
\begin{longtable}[]{@{}
  >{\raggedright\arraybackslash}p{(\linewidth - 4\tabcolsep) * \real{0.2973}}
  >{\raggedright\arraybackslash}p{(\linewidth - 4\tabcolsep) * \real{0.4324}}
  >{\raggedright\arraybackslash}p{(\linewidth - 4\tabcolsep) * \real{0.2703}}@{}}
\toprule\noalign{}
\begin{minipage}[b]{\linewidth}\raggedright
Mechanism
\end{minipage} & \begin{minipage}[b]{\linewidth}\raggedright
Formal Notation
\end{minipage} & \begin{minipage}[b]{\linewidth}\raggedright
Function
\end{minipage} \\
\midrule\noalign{}
\endhead
\textbf{Cognitive Firewall} &
\(\mathcal{F}(m) \to \{\text{accept}, \text{quarantine}, \text{reject}\}\)
& Classifies incoming messages by trust score against threshold
\(\tau\); quarantines ambiguous inputs for sandboxed evaluation \\
\textbf{Belief Sandboxing} &
\(\mathcal{B}_{\text{verified}} / \mathcal{B}_{\text{provisional}}\)
partition & Isolates unverified beliefs from the agent's operational
belief set; requires corroboration for promotion \\
\textbf{Behavioral Invariants} & Predicates
\(\text{INV}_k(\sigma_i) \in \{\text{true}, \text{false}\}\) &
Pre-defined runtime checks (e.g., goal alignment, permission boundaries)
that trigger alerts on violation \\
\textbf{Drift Detection} &
\(S_{\text{drift}} = \KL(\mathcal{B}_i^t \| \mathcal{B}_i^{t-1})\) &
Monitors belief distribution changes via KL divergence; flags sudden
shifts exceeding threshold \(\epsilon\) \\
\textbf{Byzantine Consensus} & \(\mathcal{B}_{\text{consensus}}\) with
quorum \(q\), requiring \(n \geq 3f+1\) & Multi-agent agreement protocol
tolerating up to \(f\) compromised agents among \(n\) total \\
\bottomrule\noalign{}
\end{longtable}
}

These mechanisms compose in series and parallel to achieve layered
defense. Paper 2 \cite{friedman2026cogsec2} demonstrates that the
recommended defense stack achieves 94--100\% detection at the parametric
design ceiling across 950 attack scenarios and four production
multiagent architectures. Paper 3 \cite{friedman2026cogsec3} translates
these results into deployment guidance, monitoring playbooks, and
cost--benefit frameworks; readers seeking engineering guidance on
instantiating the mechanisms below in production should consult Part 3.
\end{frame}

\begin{frame}[allowframebreaks]{Adversary Classification}
\protect\phantomsection\label{adversary-classification}
We adopt the five-tier adversary taxonomy from Paper 1
\cite{friedman2026cogsec1}:

{\def\LTcaptype{none} % do not increment counter
\begin{longtable}[]{@{}
  >{\raggedright\arraybackslash}p{(\linewidth - 6\tabcolsep) * \real{0.1628}}
  >{\raggedright\arraybackslash}p{(\linewidth - 6\tabcolsep) * \real{0.1628}}
  >{\raggedright\arraybackslash}p{(\linewidth - 6\tabcolsep) * \real{0.3023}}
  >{\raggedright\arraybackslash}p{(\linewidth - 6\tabcolsep) * \real{0.3721}}@{}}
\toprule\noalign{}
\begin{minipage}[b]{\linewidth}\raggedright
Class
\end{minipage} & \begin{minipage}[b]{\linewidth}\raggedright
Scope
\end{minipage} & \begin{minipage}[b]{\linewidth}\raggedright
Access Level
\end{minipage} & \begin{minipage}[b]{\linewidth}\raggedright
Example Vector
\end{minipage} \\
\midrule\noalign{}
\endhead
\(\Omega_1\) & External & User input & Prompt manipulation, jailbreak
attempts \\
\(\Omega_2\) & Peripheral & Tool/API, data channels & Data poisoning,
malicious web content, sensor spoofing \\
\(\Omega_3\) & Agent-level & Single compromised agent & Goal hijacking
of individual agent, compromised subagent \\
\(\Omega_4\) & Coordination & Inter-agent channels & Trust manipulation,
man-in-the-middle on agent communication \\
\(\Omega_5\) & Systemic & Orchestrator & Framework compromise, training
data manipulation \\
\bottomrule\noalign{}
\end{longtable}
}

A key observation across all ten domains analyzed in this paper is that
the primary attack vector is consistently \(\Omega_2\) (Peripheral):
adversaries inject malicious content through data channels---sensor
readings, API responses, log files, market data, diplomatic
cables---rather than through direct prompt manipulation. This reflects
the operational reality that deployed multiagent systems in critical
domains typically have hardened \(\Omega_1\) boundaries but remain
vulnerable at the data ingestion layer. Notably, \(\Omega_4\)
coordination attacks have been demonstrated empirically by He et
al.~\cite{he2025redteaming}, who introduced the Agent-in-the-Middle
(AiTM) attack exploiting inter-agent communication channels---a threat
class that our \(\Omega_2\)-focused domain analysis acknowledges but
does not examine in depth (see Limitations,
\cref{sec:limitations_discussion}).
\end{frame}

\begin{frame}[allowframebreaks]{Domain Analysis Template}
\protect\phantomsection\label{domain-analysis-template}
Each domain in this paper is analyzed using a standardized five-step
procedure:

\textbf{(i) Operational Characterization.} Identify the autonomous
agents operating in the domain, their Functional Requirements
(\(\text{FR}_k\)), Design Parameters (\(\text{DP}_k\)), and the baseline
(uncoupled) Design Matrix \([A]\).

\textbf{(ii) Attack Surface Analysis.} Describe the Goal Hijacking
attack vector, classify it by adversary class \(\Omega_k\), and identify
which OODA phase is targeted. Characterize the attack as one of three
universal patterns: FR Polarity Inversion, Constraint Relaxation, or
Context Boundary Violation.

\textbf{(iii) Transient Coupling Analysis.} Show how the attack
transforms the Design Matrix from uncoupled \([A]\) to coupled \([A']\)
by introducing off-diagonal terms. Demonstrate that the coupling is
\emph{transient}---induced by a fast OODA signal rather than a
persistent structural flaw.

\textbf{(iv) Defense Mapping.} Map the domain-specific defense strategy
to one or more of the five canonical CIF mechanisms. Where
domain-specific instantiations introduce genuinely novel patterns (e.g.,
physics-informed invariants, verification channel separation), these are
identified and elevated as contributions.

\textbf{(v) Validation Anchoring.} Cross-reference the defense mapping
to Paper 2's benchmark results \cite{friedman2026cogsec2}, confirming
that the proposed CIF mechanisms have demonstrated efficacy against the
relevant adversary class. Where appropriate, we additionally anchor
deployment-level considerations to Paper 3's operator posture and
incident-response guidance \cite{friedman2026cogsec3}.
\end{frame}

\begin{frame}[allowframebreaks]{Domain Selection Criteria}
\protect\phantomsection\label{domain-selection-criteria}
The ten domains were selected to ensure comprehensive coverage along
three dimensions:

\begin{enumerate}
\item
  \textbf{Adversary class coverage.} While all domains share
  \(\Omega_2\) as the primary vector, the specific data channels
  exploited span the full range: sensor data (Drones, Infrastructure),
  API metadata (Supply Chain), market signals (Food Security, Trade
  Wars), diplomatic communications (Nation-State), log files
  (Cyber-Security), research documents (Biowarfare), geological surveys
  (Rare Earth), and media content (Fake News).
\item
  \textbf{Temporal scale diversity.} The OODA loop operates at vastly
  different time scales across domains: milliseconds (Drone swarm
  consensus), seconds (Cyber-security incident response), hours (Supply
  chain routing), days (Infrastructure grid management), weeks (Trade
  policy), months (Nation-state diplomacy), and years (Rare Earth mining
  operations). This diversity stress-tests CIF's temporal assumptions.
\item
  \textbf{CIF mechanism coverage.} All five canonical CIF mechanisms
  appear across the domain portfolio, with each mechanism serving as the
  primary defense in at least two domains (see Discussion,
  \cref{sec:mechanism_coverage}).
\end{enumerate}
\end{frame}

\begin{frame}[allowframebreaks]{Scope and Assumptions}
\protect\phantomsection\label{scope-and-assumptions}
This analysis is qualitative and scenario-based. Each domain presents a
representative attack scenario constructed from documented real-world
incidents and known vulnerability classes, analyzed through the
CIF-AD-OODA lens. We do not claim empirical validation of defense
effectiveness in specific deployments---that requires domain-specific
experimentation beyond the scope of a cross-domain survey. However, our
methodology is informed by several concurrent frameworks: the MAESTRO
framework \cite{csa2025maestro} provides complementary layered threat
modeling for multi-agent architectures, NIST AI 600-1
\cite{nist2024genai} establishes the GenAI-specific risk profile that
our domain analysis extends to multiagent contexts, and the Agent
Security Bench (ASB) \cite{zhang2025asb}---the most comprehensive
benchmark to date with 10 scenarios, 400+ tools, and 27 attack/defense
methods---provides empirical grounding for the defense gap (84.3\%
attack success vs.~19.7\% defense success) that CIF's architectural
approach seeks to close.

We assume that the OODA loop, while a simplification of real-world
decision processes, provides a useful abstraction for identifying the
temporal dynamics of Goal Hijacking. More complex decision architectures
(e.g., nested OODA loops, parallel processing streams) would require
extensions to this model.

Finally, we assume that domain-specific latency constraints---the time
available for defense mechanisms to intervene before an attack
completes---are addressable through appropriate engineering of CIF
mechanism parameters (\(\tau\), \(\epsilon\), \(q\), \(\Delta t\)). The
determination of optimal parameter values for each domain is left for
future work.
\end{frame}

\begin{frame}[allowframebreaks]{Application to Critical Domains}
\protect\phantomsection\label{application-to-critical-domains}
The following ten sections apply this five-step CIF-AD-OODA template to
domains spanning resource extraction, geopolitics, cybersecurity,
autonomous warfare, supply chains, biosecurity, food systems, trade
policy, critical infrastructure, and information integrity. Each domain
instantiates the framework against domain-specific adversary models,
attack vectors, and defense compositions, demonstrating both the
generality of the CIF-AD-OODA integration and the domain-specific
adaptations required for effective cognitive security.
\end{frame}

\end{document}
